\section{Introduction}

Hospital reimbursement systems worldwide rely on Diagnosis-Related Groups (DRGs) to classify inpatient stays into clinically coherent categories with similar resource consumption profiles \cite{fetter1980drg}. In Chile, the public healthcare insurer FONASA adopted the IR-GRD v3.1 grouper to standardize billing across its network of public hospitals, linking reimbursement rates to the assigned DRG for each patient stay \cite{fonasa2016}. Accurate DRG assignment is thus a critical operational and financial concern for public hospitals.

The DRG grouping process maps a combination of clinical inputs---primary and secondary diagnoses coded in ICD-10 (CIE-10), procedures coded in ICD-9-PCS (CIE-9), patient age, and sex---to a single DRG class. Traditionally, this mapping is performed by rule-based grouper software that applies decision trees defined by the grouper vendor. However, the quality of DRG assignment depends heavily on the completeness and accuracy of the input codes, which are entered by clinical coders under time pressure and are susceptible to errors and omissions \cite{seroussi2017drg}.

Machine learning offers an alternative approach to DRG assignment. Rather than relying on hand-crafted decision rules, supervised ML models can learn the mapping from clinical inputs to DRG classes directly from historical data, potentially capturing patterns that rule-based systems miss. Several studies have demonstrated the feasibility of ML-based DRG prediction using structured clinical codes \cite{gartner2015ml, rajkomar2018scalable}, though most work has focused on US or European healthcare systems using CMS-DRG or G-DRG groupers.

The Chilean public health context presents a distinct setting. Hospital El Pino, located in the southern metropolitan area of Santiago, serves a predominantly low-income urban population and generates a volume of approximately 14,500 inpatient stays in the study period. The hospital's clinical coding practices, case mix, and the IR-GRD v3.1 grouper with FONASA 2016 pricing constitute a combination that has not been previously studied in the ML-for-DRG literature.

This paper presents the methodological groundwork for developing a DRG predictor for Hospital El Pino. Specifically, we (1)~describe the dataset and its characteristics, (2)~propose a preprocessing pipeline for the unordered sets of ICD codes, (3)~identify candidate ML techniques and justify their selection, (4)~define evaluation metrics appropriate for the multi-class imbalanced classification setting, and (5)~present an exploratory data analysis that characterizes the data quality, variable distributions, and class structure of the target DRG variable.
